\zlabel{triptych:concise:commentary:start}
\section{The Propers: Detailed Commentary}

\subsection{Nearer than when we first believed}

Chrysostom asks what Paul would say to someone who can remember being more eager to serve the Lord when he first believed. Paul's claim that salvation is nearer presents a difficulty for such a man. As time goes by he may actually find it harder to rouse himself. Chrysostom asks him to think of runners nearing the end of a race. The closer they come to the prize, the harder they run. Paul tells the Romans what is near enough to call for such an effort. He calls it their salvation. Chrysostom knows the day of judgment can be fearful for a man who has neglected what Christ commands. Here Paul encourages him to prepare by telling him what he may still receive.\footnote{John Chrysostom, \work{Homilies on Romans}, Homily 24, on Rom 13:11. Chrysostom and Thomas comment on the biblical introduction about knowing the time. The Latin lesson at Mass begins with an address to the brethren and retains the hour for waking. Both texts contain the comparison with the time of first believing.}

Thomas Aquinas asks what sort of distance a Christian can cover. The day on which he will answer for his life draws nearer, whatever use he makes of his years. But two people can stand at the door of a wedding and only one be ready to go in when the bridegroom arrives. Thomas recalls the people admitted to the wedding in Matthew to explain a second way of drawing near. The person who has begun to believe can grow in the charity with which he acts. He is better prepared for eternal life because of what he has become. Chrysostom's runners show why the approach of that day should stir a Christian to greater effort. Thomas explains how loving his neighbor better today prepares him for the one he hopes to meet.\footnote{Thomas Aquinas, \work{Super Romanos}, cap. 13, lectio 3, on Rom 13:11, citing Mt 25:10.}

At Mass in Advent, Christians apply the same comparison to a birth that took place long before any of them believed. Thomas asks in whose person the Church can then be speaking. Think of the people who first believed that the Son would come in the flesh. Long after they died, others who had received their promise could look back over the years in which it had been handed down. As prophecies of his coming grew more frequent, those later believers could speak of a day nearer to them than to their forebears. The Church takes up their expectation when she reads Paul in Advent. Thomas first explains the apostle's exhortation to his own contemporaries, who already knew of the Lord's birth. He distinguishes it from this use the Church makes of his words.\footnote{Aquinas, \work{Super Romanos} 13, lectio 3, [86224], paragraph on the Church's use of Paul's words in Advent.}

In the Introit one person speaks and then appeals to the experience of everyone who waits on God. The Gradual repeats just that promise and the prayer for instruction. A Christian can sing them with those who once awaited the birth, even as he prepares for another coming altogether. Just before the lesson begins, Paul speaks of the love we owe our neighbor. One man's neighbor may need the time he has been spending on an old dispute. Giving him that time would be one way of paying the debt. Neither Christmas nor the last day need first arrive for someone to take that step.

\subsection{What could bring the singer to shame?}

A man who has just told God that he trusts him asks to be spared shame. In Bellarmine's account, his enemies watch to see what will happen to someone who has staked his hopes on God. If he is left alone in his danger, they will have their occasion to laugh. In Augustine's account a person can suffer still greater harm at the hands of enemies who seem to admire him. They can congratulate him on conduct of which he ought to be ashamed. Their flattery encourages him to continue until he has disgraced himself by his own acts. When Augustine hears Christ speak the psalm in the person of the Church, it is a people acquainted with that kind of shame who say they have lifted their soul to God. They can remember what became of them when they trusted themselves.\footnote{Ps 24(25):1--4. Bellarmine, \work{Commentary on the Book of Psalms}, Ps 24, on vv. 1--4. Augustine, \work{Enarrationes}, English Psalm XXV, §§1--5.}

Someone afraid of being laughed at might think safety lay in winning the approval of the people who laughed. Augustine gives him reason to be wary of the very moment he succeeds. Bellarmine can explain why the singer goes directly from his account of his enemies to a request to be shown God's ways. He wants to live in a way pleasing to God. Only God can teach him what that involves. The assembly makes the same request again after Paul has named envy among the deeds of darkness. A man can tell his friends the story of a quarrel in such a way that they will earnestly take his side. Their support may make it harder for him to acknowledge his pleasure in doing his rival an injury. In praying to know the Lord's paths he has asked for an answer he may never hear from his friends.

When this part of the psalm returns in the Gradual the singers leave some words unsung. In the Bible, immediately after the assurance about everyone who waits, the psalmist asks that those who act unjustly be put to shame. The Gradual passes straight from the assurance to the prayer for instruction. It also changes how the assurance is sung. A man who once spoke about enemies now addresses the Lord himself at the end of the response. The appeal to be taught follows in the singular. Each of the people who just heard Paul name the conduct he expects of Christians can ask God to teach it to him.\footnote{Ps 24(25):3--4. Compare the Introit verse and the Gradual in the 1962 Missal, First Sunday of Advent, nos. 1, 4. The final address in the response also appears in Laurence Shepherd's English of the Gradual in Guéranger, \work{Advent}, p. 134.}

Chrysostom describes the sound in a street where people are beginning to wake. One has heard the early bird, and calls to the others. They can still see darkness outside, but they know enough about the approaching day to get up and put on their clothes. Paul asks something comparable of the Romans. Someone who has heard him read can begin to act on that knowledge even while a quarrel is still going in his favor. For the disciples in Luke, too, the moment for lifting their heads comes while other people can see only grounds for fear. What they have learned from Jesus about their future permits a different judgment about what they see around them now.\footnote{Chrysostom, Homily 24 on Romans, on Rom 13:12--14. Lk 21:25--28.}

\subsection{What the frightened can believe about the Judge}

Jesus had been speaking of Jerusalem surrounded by armies before he spoke of distress among the nations. Within the passage read at Mass the people's fear has a further cause. They are suffering now from their expectation of what is still to come. His disciples are there to hear him speak about people who faint. Then he addresses them directly and says what they may expect for themselves. Gregory preached this Gospel to men and women who could remember a storm in which neighbors had died overnight after making plans for the next day. They knew what menaced houses and streets could look like. Gregory had still to explain how a hearer who had lost so much might actually desire to see the Lord.\footnote{Lk 21:20--28. Gregory, \work{Homiliae in Evangelia} I.1.1--3, 5.}

Gregory asks his hearers to picture someone worn out by the road who still wants the journey to go on. What has become of his desire to reach home? A survivor of the storm might long for the chance to live in his old neighborhood again. Even the loss of much that he loved cannot of itself give him a desire for God. On Gregory's account, the disciples can raise their heads because they love the one they are about to meet. When Gregory recalls the neighbors who died in the storm, he asks the survivors how prepared they are to meet the Lord. The assurance of redemption comes from Jesus’ words to the disciples.

The Lord who will appear in majesty is also the Son whom the Church awaits at Christmas. Ambrose sees a difficulty in the cloud. Ordinarily a thick cloud obscures what a person might otherwise see. Yet in Scripture God has spoken to his people out of a cloud. Ambrose has the plural, clouds, before him in the Gospel, where the Missal has one cloud. He thinks of the many people through whom God has spoken. Moses and the prophets can be understood as clouds in which people hear Christ. For the light cloud in Isaiah's prophecy Ambrose has another explanation. The angel tells the Virgin that the power of the Most High will overshadow her. The person of whom the prophets spoke can be borne by Mary and made visible in the flesh. These figures concern how the Son has made himself known. Luke also promises that people will see him coming in power.\footnote{Ambrose, \work{Expositio evangelii secundum Lucam} X.39--43, particularly X.41--42, on Lk 21:27, Is 19:1 and Lk 1:35.}

One person whose voice is heard before the Gospel is the psalmist asking to see mercy. For Augustine the salvation he asks for is Christ himself. Those who sing the Alleluia awaiting the Nativity may want to see the face of Mary's Son. Those who sing it awaiting the last judgment may find it hard even to want that meeting. They admitted in the Collect what their own sins have done. Yet it was to the Son himself that they turned for protection. A person who can ask his Judge for help before that day already has a reason to look for him.

\subsection{What can be known from the trees}

Jesus has described events in the heavens that none of his hearers could predict. Then he speaks to people who know their trees well enough to tell when summer is near. For Gregory, the fruit this world bears is its own ruin. For someone who longs to see God, the signs also promise summer. Gregory pictures the meeting with God as summer light after the clouds of sorrow have cleared. He asks his hearers to consider whether that is the home they hope to reach.\footnote{Gregory, \work{Homiliae in Evangelia} I.1.3, on Lk 21:29--31.}

Ambrose compares Luke with Matthew. Someone remembering only Matthew's saying might picture a single fig tree. Luke also mentions all the trees. Ambrose imagines the different peoples who have come to confess God, including the Jewish people. The fruit they bear in confessing him promises a harvest for them all in the resurrection. On this account, a believer can look to the faith of people among whom the Gospel has taken root for grounds to hope for that harvest. In Gregory he was considering what gives way as the kingdom approaches. In Ambrose he has something to look for in the people who will rise.\footnote{Ambrose, \work{Expositio evangelii secundum Lucam} X.44--45, comparing Lk 21:29--30 and Mt 24:32.}

Ambrose also sees a danger in a flourishing fig tree. Elsewhere in Luke a gardener is given time to tend a tree on which the owner has found no fruit. In Genesis it is fig leaves that people take to cover their shame. Someone might be growing more confident in doing wrong and seem, by every outward test of prosperity, to have every reason to be pleased with himself. For Ambrose such growth can portend judgment. The dark interpretation comes from looking at leaves in the light of those other passages. When the words of the Gospel are read at this Mass, the trees Jesus points to are bearing fruit.\footnote{Ambrose, \work{Expositio} X.44--45, with Lk 13:6--9 and Gen 3:7.}

How much time do the signs allow? Jesus speaks of the generation that will remain until his words are fulfilled. Bede understands the generation as humankind, or as the Jewish people more particularly. Either could be spoken of across many lifetimes. Bede dwells on the solemn Amen with which the Lord gives his assurance. Truth itself has chosen to speak with the force of an oath so that even a sleepy hearer may attend. Gregory quotes the saying and follows the next verse to a different difficulty. Spoken words disappear almost before the person hearing them can answer. The earth beneath his feet seems quite solid. Yet the judgment declared by Christ in words will stand when heaven and earth have passed away. When a disciple recognizes the approach of the kingdom, he can trust what he heard Christ say about his redemption.\footnote{Bede, \work{In Lucae evangelium expositio} VI, at Lk 21:32, PL 92, col. 590B--C. Gregory, \work{Homiliae in Evangelia} I.1.4, on Lk 21:32--33.}

\subsection{The fruit of our earth}

In the Communion chant it is the earth that gives its fruit, and the verb for giving has already occurred once before the earth is even mentioned. The Lord himself will give goodness. The same people who first stood before God as suppliants at the entrance can sing of something that will come from their own ground. Does the psalm speak of a change in those who receive God's goodness, or does it tell them whom they will receive? Augustine and Bellarmine each begin one verse earlier, where truth has sprung from earth. At that verse both can speak of Mary bearing Christ.\footnote{Ps 84:12--13 (85:11--12). Augustine, \work{Enarrationes}, English Psalm LXXXV, §§10--12. Bellarmine, \work{Commentary on the Book of Psalms}, Ps 84, exposition 12, p. 259.}

For Augustine it is astonishing that truth should spring from the very earth of which human beings are made. They are capable of lying even about their own conduct. Yet the one who is Truth has received from a woman of this earth the mortal flesh in which he can offer himself. Augustine finds a further instance of truth springing from earth. A person who confesses his sin has begun to say something true about himself. That image of confession allows Augustine to picture an earth already ploughed by the time he comes to the words sung at Communion. His hearers have opened furrows by their confession. Hearing God's word, they have received the seed. A field ploughed and sown still needs rain.

The Missal promises goodness. In the English translation of Augustine's sermon the Lord gives sweetness. A man who once took pleasure in stealing another person's property now takes such pleasure in giving his own away that he goes looking for someone to share it with. A man capable of enjoying such an act can begin to understand Gregory's claim that the disciples may look forward to meeting their Judge. He can desire to meet the Lord who has given him such delight in another person's good.\footnote{Augustine, \work{Enarrationes}, English Psalm LXXXV, §§11--12. Compare the Communion of the 1962 Missal, First Sunday of Advent, no. 9.}

Bellarmine wants to account for a field that bears fruit without anyone having cultivated it. He asks how the Son who comes from God can spring from the earth, as the preceding verse says. Picture a flower growing wild. Nobody has brought the plough there or planted seed in that spot. Yet the earth really bears the flower, given the rain and the light it receives. Where the psalm speaks of what the Lord gives, Bellarmine sees the work of the Spirit overshadowing Mary. Where the earth gives its fruit, he sees the Son born of her. To call it our earth is to confess that the Son receives his flesh from one of our own. Augustine needed his hearers to recognize ground cultivated in confession and still awaiting rain. For Bellarmine it is precisely the absence of human cultivation that lets him speak of a virginal birth.\footnote{Bellarmine, Ps 84, exposition 12, O'Sullivan's abridged English, p. 259. The Communion is Douay Ps 84:13.}

Bellarmine explains why a man approaching Communion might think of the child whom Mary bore. Augustine lets a man who knows he once stole from his neighbor consider what the Lord's goodness may yet produce in him. At the Secret the priest asked that the gifts might cleanse those who come to their giver. Such a man can have both a person to approach and a particular gift to ask of him.

\subsection{After the people have received}

At the Offertory the Church brings bread and wine while keeping her gaze on the one she expects. In the Latin she sings again the very antiphon used at the entrance. At the entrance the psalm verse and the doxology followed it. Here at the altar the people have the first-person declaration of trust again without those additions. The danger the singer named at the entrance still has a place in his prayer after the Gospel. Even as he offers something of his own, the suppliant tells God upon whom he relies.\footnote{Guéranger, \work{The Liturgical Year}, Advent, First Sunday of Advent, pp. 135--136. Augustine, \work{Enarrationes}, English Psalm XXV, §§1--5. 1962 Missal, First Sunday of Advent, nos. 1, 7.}

The Collect was addressed to the Son. In the Secret the priest prays to the Father through him, asking that the worshippers may approach the source of the gifts. In the historical English of Lasance and Walsh that source is called their author. The thanksgiving that follows declares the Son and the Spirit equal to the Father in majesty. A person might find it hard to reconcile that equality with the thought of the Son being sent into the world. Augustine, in his work on the Trinity, understands being sent by reference to the relation of the Son to the Father who sends him. The Son who took flesh from Mary retains the divine majesty the faithful adore when they give thanks.\footnote{1962 Missal, First Sunday of Advent, nos. 2, 8, and Preface of the Most Holy Trinity, p. 293, no. 1082. Lasance and Walsh, \work{The New Roman Missal}, pp. 103, 105. Augustine, \work{De Trinitate} III, preface §3.}

Guéranger describes Mary made fruitful by the dew of heaven in his comment on the Communion. When the priest and people have received Communion, he says, they can still sing of what the earth will give. They can expect with those who waited for the birth because the Son has taken flesh. And because he gives himself to them in Communion they already know him as the one they expect. His gift also lets them look beyond the feast. In the Gospel the Son whom his mother bore will be visible to men in his majesty. The believer who has learned to recognize him under sacramental signs has still to see him face to face.\footnote{Guéranger, \work{Advent}, First Sunday of Advent, p. 137, before the Communion antiphon.}

The Church's last petition gives the communicant something to ask for in the meantime. The psalmist remembered receiving mercy in the temple. Now the priest asks that he and the people may receive it there so that they may prepare to celebrate their restoration. Even the person who has received Communion a few moments before can still need help to keep the feast. If he has been waiting for the Judge, he can ask how to give up the quarrel Paul forbids. Chrysostom compares putting on Christ to the way people speak of two friends who are always together. Christ, he says, can be seen in us when we do what he did. Someone who has been hurt in a quarrel might be the first to notice an unexpected gentleness in his adversary. If he has been awaiting the birth, he can ask how to honor a gift that came to him before he could do anything to deserve it. Thomas's distinction between elapsed time and growth in charity leaves room for an answer before the feast arrives. The neighbor to whom Paul says he owes love may already be at hand.\footnote{Rom 13:8, 13--14. Chrysostom, Homily 24 on Romans, on putting on Christ. Aquinas, \work{Super Romanos} 13, lectio 3. 1962 Missal, First Sunday of Advent, no. 10.}
